% Paper A: selection-robust violation guarantee.
% Included by paper.tex (% Paper A: selection-robust violation guarantee.
% Included by paper.tex (% Paper A: selection-robust violation guarantee.
% Included by paper.tex (% Paper A: selection-robust violation guarantee.
% Included by paper.tex (\input{paper_a_guarantee}) in Section 3. Uses the
% assumption, definition, theorem, corollary and remark environments and the
% \status macro defined there. Citation keys are those of paper_a.bib.
%
% Proofs other than Theorem 1's are in paper_a_proofs.tex (the appendix).
% Results held until probe P1 decides the chemotherapy design (joint radius,
% certified level, simulated level) are in paper_a_guarantee_held.tex, which
% is not \input.
%
% \status tags. JS checked by hand the version committed in 40acea5: Theorem 1,
% Corollaries 1-3 and the closing remark, for Sigma = sigma^2 I. The general
% Sigma, the remark on the choice of norm and everything from "The simulated
% radius" on came later and are not checked.

\begin{outline}
Tags: JS checked by hand the $\Sigma = \sigma^2 I$ version of
Theorem~\ref{thm:selection}, Corollaries~\ref{cor:finite}--\ref{cor:multi}
and the closing remark (commit \texttt{40acea5}). The general $\Sigma$, the
remark on the choice of norm and everything after the closing remark were
added later and are not checked.
\end{outline}

\paragraph{Setting.}
We observe training data $(x_i, y_i)_{i=1}^n$ with $y = y^\star + \varepsilon$,
where $y^\star_i = h^\star(x_i)$ are the noise-free labels. The learned
constraint is $f(x;\theta(y)) \le b$ for $x \in \mathcal{X}$, where $\theta(y)$
denotes the model fitted on labels $y$.

\begin{assumption}[Label noise]\label{ass:noise}
\status{checked by JS for $\Sigma = \sigma^2 I$; general $\Sigma$ not checked}
The noise is $\varepsilon = \Sigma^{1/2}\xi$, where $\Sigma \succeq 0$ is known
and $\xi_1,\dots,\xi_n$ are independent, zero-mean and sub-Gaussian with
variance proxy $1$. We write $\|v\|_\Sigma = \|\Sigma^{1/2}v\|_2$.
\end{assumption}

With $\Sigma = \sigma^2 I$ the noise terms are independent with variance proxy
$\sigma^2$ and $\|v\|_\Sigma = \sigma\|v\|_2$. A diagonal $\Sigma$ allows
unequal noise levels, and a full one correlated noise.

\begin{assumption}[Affine predictor]\label{ass:affine}
For all $x \in \mathcal{X}$, $f(x;\theta(y)) = s(x)^\top y + c(x)$, where
$s(x) \in \mathbb{R}^n$ and $c(x)$ do not depend on $\varepsilon$.
\end{assumption}

\begin{definition}[Clean-label model]\label{def:clean}
Let $\hat\theta = \theta(y)$ and $\theta^\star = \theta(y^\star)$. The
clean-label model is $f(\cdot\,;\theta^\star)$.
\end{definition}

Since $s(x)$ and $c(x)$ do not depend on $\varepsilon$,
Assumption~\ref{ass:affine} applies to $y^\star = y - \varepsilon$ as well as to
$y$, and subtracting the two gives, for all $x \in \mathcal{X}$,
\begin{equation}\label{eq:identity}
  f(x;\theta^\star) = f(x;\hat\theta) - s(x)^\top \varepsilon .
\end{equation}

For a radius $R > 0$ fixed independently of $\varepsilon$, the robust
counterpart of the learned constraint is
\begin{equation}\label{eq:rc}
  f(x;\hat\theta) + R\,\|s(x)\|_\Sigma \le b .
\end{equation}
The radius $R$ is in units of the noise's scale: with $\Sigma = \sigma^2 I$ the
margin is $R\sigma\|s(x)\|_2$.

\begin{theorem}\label{thm:selection}
\status{checked by JS for $\Sigma = \sigma^2 I$; general $\Sigma$ not checked}
Suppose Assumptions~\ref{ass:noise} and~\ref{ass:affine} hold, and let $x^\star$
be any solution satisfying \eqref{eq:rc}, possibly chosen as a function of the
data. Then
\[
  \mathbb{P}\bigl(f(x^\star;\theta^\star) > b\bigr) \le \mathbb{P}(W > R),
  \qquad
  W = \sup_{x \in \mathcal{X},\, \|s(x)\|_\Sigma \neq 0}
      \frac{-s(x)^\top \varepsilon}{\|s(x)\|_\Sigma} .
\]
\end{theorem}

\begin{proof}
Suppose $f(x^\star;\theta^\star) > b$. By \eqref{eq:identity} and
\eqref{eq:rc},
\[
  -s(x^\star)^\top \varepsilon
  = f(x^\star;\theta^\star) - f(x^\star;\hat\theta)
  > b - f(x^\star;\hat\theta)
  \ge R\,\|s(x^\star)\|_\Sigma .
\]
The left side equals $-(\Sigma^{1/2}s(x^\star))^\top\xi$ and is nonzero, so
$\|s(x^\star)\|_\Sigma \neq 0$ and $W > R$.
\end{proof}

\begin{remark}[Choice of norm]\label{rem:norm}
\status{not checked: added after JS's check}
The proof uses only the inequality
$-s(x^\star)^\top\varepsilon > R\,\|s(x^\star)\|_\Sigma$, so
Theorem~\ref{thm:selection} holds verbatim for any norm $N$ on $\mathbb{R}^n$
in place of $\|\cdot\|_\Sigma$, both in \eqref{eq:rc} and in $W$. The norm
enters only through the distribution of $W$. Under $\|\cdot\|_\Sigma$ the
standardized error at $x$ is
$-s(x)^\top\varepsilon/\|s(x)\|_\Sigma = -u(x)^\top\xi$ with $u(x)$ a unit
vector, so it has the same tail at every $x$; Corollaries~\ref{cor:finite}
and~\ref{cor:subspace} use exactly this. Under another norm the standardized
error at $x$ is that one multiplied by
$\rho(x) = \|s(x)\|_\Sigma / N(s(x))$. Bounding the corresponding $W_N$
through $W$, as $W_N \le \bar\rho\max(W,0)$ with
$\bar\rho = \sup_{x\in\mathcal{X}} \rho(x)$, turns a radius $R$ from these
corollaries into $\bar\rho R$, and the margin
$\bar\rho R\,N(s(x)) \ge R\,\|s(x)\|_\Sigma$ is at least as large at every
$x$. For example, noise bounded as $|\varepsilon_i| \le c$ gives the box
counterpart $f(x;\hat\theta) + c\,\|s(x)\|_1 \le b$, which holds with
probability one; but $\|s(x)\|_1$ can be up to $\sqrt{n}$ times
$\|s(x)\|_2$.
\end{remark}

\begin{corollary}[Tree ensembles]\label{cor:finite}
\status{checked by JS for $\Sigma = \sigma^2 I$; general $\Sigma$ not checked}
If $s(x)$ takes at most $K$ distinct values on $\mathcal{X}$, then
$R = \sqrt{2\ln(K/\epsilon)}$ guarantees
$\mathbb{P}\bigl(f(x^\star;\theta^\star) > b\bigr) \le \epsilon$.
\end{corollary}

\begin{corollary}[Linear models]\label{cor:subspace}
\status{checked by JS for $\Sigma = \sigma^2 I$; general $\Sigma$ not checked}
Suppose $s(x)$ lies in a subspace $S \subseteq \mathbb{R}^n$ of dimension $d$
for all $x \in \mathcal{X}$, and $\xi \sim N(0, I)$, that is,
$\varepsilon \sim N(0,\Sigma)$. Then $R = \sqrt{\chi^2_{d,1-\epsilon}}$
guarantees $\mathbb{P}\bigl(f(x^\star;\theta^\star) > b\bigr) \le \epsilon$,
where $\chi^2_{d,1-\epsilon}$ is the $(1-\epsilon)$-quantile of the chi-square
distribution with $d$ degrees of freedom.
\end{corollary}

\begin{corollary}[Multiple learned constraints]\label{cor:multi}
\status{checked by JS for $\Sigma = \sigma^2 I$; general $\Sigma$ not checked}
Consider $m$ learned constraints, each satisfying
Assumptions~\ref{ass:noise} and~\ref{ass:affine} with its own $s_j$,
$\Sigma_j$ and radius $R_j$. If each $R_j$ is chosen as in
Corollary~\ref{cor:finite} or~\ref{cor:subspace} at level $\epsilon/m$, then
the probability that $x^\star$ violates at least one clean-label constraint is
at most $\epsilon$. No assumption is required on the dependence of the noise
across constraints.
\end{corollary}

\begin{remark}
\status{checked by JS for $\Sigma = \sigma^2 I$}
(i) For ridge regression, $s(x) = Z(Z^\top Z + \lambda I)^{-1}\phi(x)$, so $S$
is the column space of the training matrix $Z$ and $d = \mathrm{rank}(Z)$.
(ii) A lower-bound constraint $f(x;\theta) \ge b$ is handled symmetrically,
with robust counterpart $f(x;\hat\theta) - R\|s(x)\|_\Sigma \ge b$ and the same
radii.
(iii) The guarantee concerns the propagation of label noise only; the
approximation error $h^\star - f(\cdot\,;\theta^\star)$ is not covered.
(iv) For models whose structure is selected from data (active sets, splits,
hidden layers), Assumption~\ref{ass:affine} requires the structure to be
chosen independently of $\varepsilon$, for instance by sample splitting.
\end{remark}

\paragraph{The simulated radius.}
Corollaries~\ref{cor:finite} and~\ref{cor:subspace} bound $W$ without using
$\mathcal{X}$: they protect against every direction among the vectors
$\Sigma^{1/2}s(x)$, whether or not a decision in $\mathcal{X}$ reaches it. When
the distribution of $\xi$ is known, the radius can instead be read off
simulated copies of $W$, and an order statistic of the copies is valid with no
sampling-error term. Write $W = w(\xi)$, where
\[
  w(z) = \sup_{x \in \mathcal{X},\, \|s(x)\|_\Sigma \neq 0}
    \frac{-(\Sigma^{1/2}s(x))^\top z}{\|s(x)\|_\Sigma},
  \qquad z \in \mathbb{R}^n ,
\]
assumed measurable. Under Assumption~\ref{ass:affine}, $s(\cdot)$ does not
depend on $\varepsilon$, so $w$ is a fixed function, known once the model's
structure is. Evaluating $w(z)$ is a maximization over $\mathcal{X}$, a
mixed-integer program when $\mathcal{X}$ is mixed-integer.

\begin{proposition}[Simulated radius]\label{prop:simulated}
\status{not checked: drafted 2026-11-15}
Suppose Assumptions~\ref{ass:noise} and~\ref{ass:affine} hold and the
distribution of $\xi$ is known. Draw $\xi^{(1)},\dots,\xi^{(N)}$ from it,
independently of each other and of $\varepsilon$, and let
$\overline W_1,\dots,\overline W_N$ be any numbers with
$\overline W_i \ge w(\xi^{(i)})$. For $\epsilon \in (0,1)$ and
$N \ge 1/\epsilon - 1$, let $k = \lceil (N+1)(1-\epsilon) \rceil$ and let
$\hat R$ be the $k$-th smallest of $\overline W_1,\dots,\overline W_N$. Then
any $x^\star$ satisfying \eqref{eq:rc} at $R = \hat R$, possibly chosen as a
function of the data and of the simulation, satisfies
$\mathbb{P}\bigl(f(x^\star;\theta^\star) > b\bigr) \le \epsilon$, where the
probability is over the noise and the simulation jointly.
\end{proposition}

\begin{remark}[Computing the radius]\label{rem:simulated}
\status{not checked: drafted 2026-11-15}
(i) Each $\overline W_i$ must be an upper bound on $w(\xi^{(i)})$. The value
of a feasible point of the maximization, or the largest standardized error
over a sample of decisions, is a lower bound, and an order statistic of lower
bounds can fall below the $(1-\epsilon)$-quantile of $W$, which makes the
radius too small. A solver's dual bound is an upper bound at any time limit,
so stopping a maximization early makes the radius larger, never invalid.
(ii) The proposition needs the distribution of $\xi$, not only a variance
proxy: the argument requires the simulated copies to have the law of the
true $W$. The experiments use $\xi \sim N(0, I)$. An estimated $\Sigma$ is
not covered yet. \tbd{T2}
\end{remark}

Corollary~\ref{cor:multi} holds with unequal levels $\epsilon_j$ summing to
$\epsilon$, and with any radius valid at level $\epsilon_j$ for constraint $j$
alone, including the simulated radius of
Proposition~\ref{prop:simulated}: its proof uses only
$\mathbb{P}(W_j > R_j) \le \epsilon_j$. \status{not checked: drafted
2026-11-15}
) in Section 3. Uses the
% assumption, definition, theorem, corollary and remark environments and the
% \status macro defined there. Citation keys are those of paper_a.bib.
%
% Proofs other than Theorem 1's are in paper_a_proofs.tex (the appendix).
% Results held until probe P1 decides the chemotherapy design (joint radius,
% certified level, simulated level) are in paper_a_guarantee_held.tex, which
% is not \input.
%
% \status tags. JS checked by hand the version committed in 40acea5: Theorem 1,
% Corollaries 1-3 and the closing remark, for Sigma = sigma^2 I. The general
% Sigma, the remark on the choice of norm and everything from "The simulated
% radius" on came later and are not checked.

\begin{outline}
Tags: JS checked by hand the $\Sigma = \sigma^2 I$ version of
Theorem~\ref{thm:selection}, Corollaries~\ref{cor:finite}--\ref{cor:multi}
and the closing remark (commit \texttt{40acea5}). The general $\Sigma$, the
remark on the choice of norm and everything after the closing remark were
added later and are not checked.
\end{outline}

\paragraph{Setting.}
We observe training data $(x_i, y_i)_{i=1}^n$ with $y = y^\star + \varepsilon$,
where $y^\star_i = h^\star(x_i)$ are the noise-free labels. The learned
constraint is $f(x;\theta(y)) \le b$ for $x \in \mathcal{X}$, where $\theta(y)$
denotes the model fitted on labels $y$.

\begin{assumption}[Label noise]\label{ass:noise}
\status{checked by JS for $\Sigma = \sigma^2 I$; general $\Sigma$ not checked}
The noise is $\varepsilon = \Sigma^{1/2}\xi$, where $\Sigma \succeq 0$ is known
and $\xi_1,\dots,\xi_n$ are independent, zero-mean and sub-Gaussian with
variance proxy $1$. We write $\|v\|_\Sigma = \|\Sigma^{1/2}v\|_2$.
\end{assumption}

With $\Sigma = \sigma^2 I$ the noise terms are independent with variance proxy
$\sigma^2$ and $\|v\|_\Sigma = \sigma\|v\|_2$. A diagonal $\Sigma$ allows
unequal noise levels, and a full one correlated noise.

\begin{assumption}[Affine predictor]\label{ass:affine}
For all $x \in \mathcal{X}$, $f(x;\theta(y)) = s(x)^\top y + c(x)$, where
$s(x) \in \mathbb{R}^n$ and $c(x)$ do not depend on $\varepsilon$.
\end{assumption}

\begin{definition}[Clean-label model]\label{def:clean}
Let $\hat\theta = \theta(y)$ and $\theta^\star = \theta(y^\star)$. The
clean-label model is $f(\cdot\,;\theta^\star)$.
\end{definition}

Since $s(x)$ and $c(x)$ do not depend on $\varepsilon$,
Assumption~\ref{ass:affine} applies to $y^\star = y - \varepsilon$ as well as to
$y$, and subtracting the two gives, for all $x \in \mathcal{X}$,
\begin{equation}\label{eq:identity}
  f(x;\theta^\star) = f(x;\hat\theta) - s(x)^\top \varepsilon .
\end{equation}

For a radius $R > 0$ fixed independently of $\varepsilon$, the robust
counterpart of the learned constraint is
\begin{equation}\label{eq:rc}
  f(x;\hat\theta) + R\,\|s(x)\|_\Sigma \le b .
\end{equation}
The radius $R$ is in units of the noise's scale: with $\Sigma = \sigma^2 I$ the
margin is $R\sigma\|s(x)\|_2$.

\begin{theorem}\label{thm:selection}
\status{checked by JS for $\Sigma = \sigma^2 I$; general $\Sigma$ not checked}
Suppose Assumptions~\ref{ass:noise} and~\ref{ass:affine} hold, and let $x^\star$
be any solution satisfying \eqref{eq:rc}, possibly chosen as a function of the
data. Then
\[
  \mathbb{P}\bigl(f(x^\star;\theta^\star) > b\bigr) \le \mathbb{P}(W > R),
  \qquad
  W = \sup_{x \in \mathcal{X},\, \|s(x)\|_\Sigma \neq 0}
      \frac{-s(x)^\top \varepsilon}{\|s(x)\|_\Sigma} .
\]
\end{theorem}

\begin{proof}
Suppose $f(x^\star;\theta^\star) > b$. By \eqref{eq:identity} and
\eqref{eq:rc},
\[
  -s(x^\star)^\top \varepsilon
  = f(x^\star;\theta^\star) - f(x^\star;\hat\theta)
  > b - f(x^\star;\hat\theta)
  \ge R\,\|s(x^\star)\|_\Sigma .
\]
The left side equals $-(\Sigma^{1/2}s(x^\star))^\top\xi$ and is nonzero, so
$\|s(x^\star)\|_\Sigma \neq 0$ and $W > R$.
\end{proof}

\begin{remark}[Choice of norm]\label{rem:norm}
\status{not checked: added after JS's check}
The proof uses only the inequality
$-s(x^\star)^\top\varepsilon > R\,\|s(x^\star)\|_\Sigma$, so
Theorem~\ref{thm:selection} holds verbatim for any norm $N$ on $\mathbb{R}^n$
in place of $\|\cdot\|_\Sigma$, both in \eqref{eq:rc} and in $W$. The norm
enters only through the distribution of $W$. Under $\|\cdot\|_\Sigma$ the
standardized error at $x$ is
$-s(x)^\top\varepsilon/\|s(x)\|_\Sigma = -u(x)^\top\xi$ with $u(x)$ a unit
vector, so it has the same tail at every $x$; Corollaries~\ref{cor:finite}
and~\ref{cor:subspace} use exactly this. Under another norm the standardized
error at $x$ is that one multiplied by
$\rho(x) = \|s(x)\|_\Sigma / N(s(x))$. Bounding the corresponding $W_N$
through $W$, as $W_N \le \bar\rho\max(W,0)$ with
$\bar\rho = \sup_{x\in\mathcal{X}} \rho(x)$, turns a radius $R$ from these
corollaries into $\bar\rho R$, and the margin
$\bar\rho R\,N(s(x)) \ge R\,\|s(x)\|_\Sigma$ is at least as large at every
$x$. For example, noise bounded as $|\varepsilon_i| \le c$ gives the box
counterpart $f(x;\hat\theta) + c\,\|s(x)\|_1 \le b$, which holds with
probability one; but $\|s(x)\|_1$ can be up to $\sqrt{n}$ times
$\|s(x)\|_2$.
\end{remark}

\begin{corollary}[Tree ensembles]\label{cor:finite}
\status{checked by JS for $\Sigma = \sigma^2 I$; general $\Sigma$ not checked}
If $s(x)$ takes at most $K$ distinct values on $\mathcal{X}$, then
$R = \sqrt{2\ln(K/\epsilon)}$ guarantees
$\mathbb{P}\bigl(f(x^\star;\theta^\star) > b\bigr) \le \epsilon$.
\end{corollary}

\begin{corollary}[Linear models]\label{cor:subspace}
\status{checked by JS for $\Sigma = \sigma^2 I$; general $\Sigma$ not checked}
Suppose $s(x)$ lies in a subspace $S \subseteq \mathbb{R}^n$ of dimension $d$
for all $x \in \mathcal{X}$, and $\xi \sim N(0, I)$, that is,
$\varepsilon \sim N(0,\Sigma)$. Then $R = \sqrt{\chi^2_{d,1-\epsilon}}$
guarantees $\mathbb{P}\bigl(f(x^\star;\theta^\star) > b\bigr) \le \epsilon$,
where $\chi^2_{d,1-\epsilon}$ is the $(1-\epsilon)$-quantile of the chi-square
distribution with $d$ degrees of freedom.
\end{corollary}

\begin{corollary}[Multiple learned constraints]\label{cor:multi}
\status{checked by JS for $\Sigma = \sigma^2 I$; general $\Sigma$ not checked}
Consider $m$ learned constraints, each satisfying
Assumptions~\ref{ass:noise} and~\ref{ass:affine} with its own $s_j$,
$\Sigma_j$ and radius $R_j$. If each $R_j$ is chosen as in
Corollary~\ref{cor:finite} or~\ref{cor:subspace} at level $\epsilon/m$, then
the probability that $x^\star$ violates at least one clean-label constraint is
at most $\epsilon$. No assumption is required on the dependence of the noise
across constraints.
\end{corollary}

\begin{remark}
\status{checked by JS for $\Sigma = \sigma^2 I$}
(i) For ridge regression, $s(x) = Z(Z^\top Z + \lambda I)^{-1}\phi(x)$, so $S$
is the column space of the training matrix $Z$ and $d = \mathrm{rank}(Z)$.
(ii) A lower-bound constraint $f(x;\theta) \ge b$ is handled symmetrically,
with robust counterpart $f(x;\hat\theta) - R\|s(x)\|_\Sigma \ge b$ and the same
radii.
(iii) The guarantee concerns the propagation of label noise only; the
approximation error $h^\star - f(\cdot\,;\theta^\star)$ is not covered.
(iv) For models whose structure is selected from data (active sets, splits,
hidden layers), Assumption~\ref{ass:affine} requires the structure to be
chosen independently of $\varepsilon$, for instance by sample splitting.
\end{remark}

\paragraph{The simulated radius.}
Corollaries~\ref{cor:finite} and~\ref{cor:subspace} bound $W$ without using
$\mathcal{X}$: they protect against every direction among the vectors
$\Sigma^{1/2}s(x)$, whether or not a decision in $\mathcal{X}$ reaches it. When
the distribution of $\xi$ is known, the radius can instead be read off
simulated copies of $W$, and an order statistic of the copies is valid with no
sampling-error term. Write $W = w(\xi)$, where
\[
  w(z) = \sup_{x \in \mathcal{X},\, \|s(x)\|_\Sigma \neq 0}
    \frac{-(\Sigma^{1/2}s(x))^\top z}{\|s(x)\|_\Sigma},
  \qquad z \in \mathbb{R}^n ,
\]
assumed measurable. Under Assumption~\ref{ass:affine}, $s(\cdot)$ does not
depend on $\varepsilon$, so $w$ is a fixed function, known once the model's
structure is. Evaluating $w(z)$ is a maximization over $\mathcal{X}$, a
mixed-integer program when $\mathcal{X}$ is mixed-integer.

\begin{proposition}[Simulated radius]\label{prop:simulated}
\status{not checked: drafted 2026-11-15}
Suppose Assumptions~\ref{ass:noise} and~\ref{ass:affine} hold and the
distribution of $\xi$ is known. Draw $\xi^{(1)},\dots,\xi^{(N)}$ from it,
independently of each other and of $\varepsilon$, and let
$\overline W_1,\dots,\overline W_N$ be any numbers with
$\overline W_i \ge w(\xi^{(i)})$. For $\epsilon \in (0,1)$ and
$N \ge 1/\epsilon - 1$, let $k = \lceil (N+1)(1-\epsilon) \rceil$ and let
$\hat R$ be the $k$-th smallest of $\overline W_1,\dots,\overline W_N$. Then
any $x^\star$ satisfying \eqref{eq:rc} at $R = \hat R$, possibly chosen as a
function of the data and of the simulation, satisfies
$\mathbb{P}\bigl(f(x^\star;\theta^\star) > b\bigr) \le \epsilon$, where the
probability is over the noise and the simulation jointly.
\end{proposition}

\begin{remark}[Computing the radius]\label{rem:simulated}
\status{not checked: drafted 2026-11-15}
(i) Each $\overline W_i$ must be an upper bound on $w(\xi^{(i)})$. The value
of a feasible point of the maximization, or the largest standardized error
over a sample of decisions, is a lower bound, and an order statistic of lower
bounds can fall below the $(1-\epsilon)$-quantile of $W$, which makes the
radius too small. A solver's dual bound is an upper bound at any time limit,
so stopping a maximization early makes the radius larger, never invalid.
(ii) The proposition needs the distribution of $\xi$, not only a variance
proxy: the argument requires the simulated copies to have the law of the
true $W$. The experiments use $\xi \sim N(0, I)$. An estimated $\Sigma$ is
not covered yet. \tbd{T2}
\end{remark}

Corollary~\ref{cor:multi} holds with unequal levels $\epsilon_j$ summing to
$\epsilon$, and with any radius valid at level $\epsilon_j$ for constraint $j$
alone, including the simulated radius of
Proposition~\ref{prop:simulated}: its proof uses only
$\mathbb{P}(W_j > R_j) \le \epsilon_j$. \status{not checked: drafted
2026-11-15}
) in Section 3. Uses the
% assumption, definition, theorem, corollary and remark environments and the
% \status macro defined there. Citation keys are those of paper_a.bib.
%
% Proofs other than Theorem 1's are in paper_a_proofs.tex (the appendix).
% Results held until probe P1 decides the chemotherapy design (joint radius,
% certified level, simulated level) are in paper_a_guarantee_held.tex, which
% is not \input.
%
% \status tags. JS checked by hand the version committed in 40acea5: Theorem 1,
% Corollaries 1-3 and the closing remark, for Sigma = sigma^2 I. The general
% Sigma, the remark on the choice of norm and everything from "The simulated
% radius" on came later and are not checked.

\begin{outline}
Tags: JS checked by hand the $\Sigma = \sigma^2 I$ version of
Theorem~\ref{thm:selection}, Corollaries~\ref{cor:finite}--\ref{cor:multi}
and the closing remark (commit \texttt{40acea5}). The general $\Sigma$, the
remark on the choice of norm and everything after the closing remark were
added later and are not checked.
\end{outline}

\paragraph{Setting.}
We observe training data $(x_i, y_i)_{i=1}^n$ with $y = y^\star + \varepsilon$,
where $y^\star_i = h^\star(x_i)$ are the noise-free labels. The learned
constraint is $f(x;\theta(y)) \le b$ for $x \in \mathcal{X}$, where $\theta(y)$
denotes the model fitted on labels $y$.

\begin{assumption}[Label noise]\label{ass:noise}
\status{checked by JS for $\Sigma = \sigma^2 I$; general $\Sigma$ not checked}
The noise is $\varepsilon = \Sigma^{1/2}\xi$, where $\Sigma \succeq 0$ is known
and $\xi_1,\dots,\xi_n$ are independent, zero-mean and sub-Gaussian with
variance proxy $1$. We write $\|v\|_\Sigma = \|\Sigma^{1/2}v\|_2$.
\end{assumption}

With $\Sigma = \sigma^2 I$ the noise terms are independent with variance proxy
$\sigma^2$ and $\|v\|_\Sigma = \sigma\|v\|_2$. A diagonal $\Sigma$ allows
unequal noise levels, and a full one correlated noise.

\begin{assumption}[Affine predictor]\label{ass:affine}
For all $x \in \mathcal{X}$, $f(x;\theta(y)) = s(x)^\top y + c(x)$, where
$s(x) \in \mathbb{R}^n$ and $c(x)$ do not depend on $\varepsilon$.
\end{assumption}

\begin{definition}[Clean-label model]\label{def:clean}
Let $\hat\theta = \theta(y)$ and $\theta^\star = \theta(y^\star)$. The
clean-label model is $f(\cdot\,;\theta^\star)$.
\end{definition}

Since $s(x)$ and $c(x)$ do not depend on $\varepsilon$,
Assumption~\ref{ass:affine} applies to $y^\star = y - \varepsilon$ as well as to
$y$, and subtracting the two gives, for all $x \in \mathcal{X}$,
\begin{equation}\label{eq:identity}
  f(x;\theta^\star) = f(x;\hat\theta) - s(x)^\top \varepsilon .
\end{equation}

For a radius $R > 0$ fixed independently of $\varepsilon$, the robust
counterpart of the learned constraint is
\begin{equation}\label{eq:rc}
  f(x;\hat\theta) + R\,\|s(x)\|_\Sigma \le b .
\end{equation}
The radius $R$ is in units of the noise's scale: with $\Sigma = \sigma^2 I$ the
margin is $R\sigma\|s(x)\|_2$.

\begin{theorem}\label{thm:selection}
\status{checked by JS for $\Sigma = \sigma^2 I$; general $\Sigma$ not checked}
Suppose Assumptions~\ref{ass:noise} and~\ref{ass:affine} hold, and let $x^\star$
be any solution satisfying \eqref{eq:rc}, possibly chosen as a function of the
data. Then
\[
  \mathbb{P}\bigl(f(x^\star;\theta^\star) > b\bigr) \le \mathbb{P}(W > R),
  \qquad
  W = \sup_{x \in \mathcal{X},\, \|s(x)\|_\Sigma \neq 0}
      \frac{-s(x)^\top \varepsilon}{\|s(x)\|_\Sigma} .
\]
\end{theorem}

\begin{proof}
Suppose $f(x^\star;\theta^\star) > b$. By \eqref{eq:identity} and
\eqref{eq:rc},
\[
  -s(x^\star)^\top \varepsilon
  = f(x^\star;\theta^\star) - f(x^\star;\hat\theta)
  > b - f(x^\star;\hat\theta)
  \ge R\,\|s(x^\star)\|_\Sigma .
\]
The left side equals $-(\Sigma^{1/2}s(x^\star))^\top\xi$ and is nonzero, so
$\|s(x^\star)\|_\Sigma \neq 0$ and $W > R$.
\end{proof}

\begin{remark}[Choice of norm]\label{rem:norm}
\status{not checked: added after JS's check}
The proof uses only the inequality
$-s(x^\star)^\top\varepsilon > R\,\|s(x^\star)\|_\Sigma$, so
Theorem~\ref{thm:selection} holds verbatim for any norm $N$ on $\mathbb{R}^n$
in place of $\|\cdot\|_\Sigma$, both in \eqref{eq:rc} and in $W$. The norm
enters only through the distribution of $W$. Under $\|\cdot\|_\Sigma$ the
standardized error at $x$ is
$-s(x)^\top\varepsilon/\|s(x)\|_\Sigma = -u(x)^\top\xi$ with $u(x)$ a unit
vector, so it has the same tail at every $x$; Corollaries~\ref{cor:finite}
and~\ref{cor:subspace} use exactly this. Under another norm the standardized
error at $x$ is that one multiplied by
$\rho(x) = \|s(x)\|_\Sigma / N(s(x))$. Bounding the corresponding $W_N$
through $W$, as $W_N \le \bar\rho\max(W,0)$ with
$\bar\rho = \sup_{x\in\mathcal{X}} \rho(x)$, turns a radius $R$ from these
corollaries into $\bar\rho R$, and the margin
$\bar\rho R\,N(s(x)) \ge R\,\|s(x)\|_\Sigma$ is at least as large at every
$x$. For example, noise bounded as $|\varepsilon_i| \le c$ gives the box
counterpart $f(x;\hat\theta) + c\,\|s(x)\|_1 \le b$, which holds with
probability one; but $\|s(x)\|_1$ can be up to $\sqrt{n}$ times
$\|s(x)\|_2$.
\end{remark}

\begin{corollary}[Tree ensembles]\label{cor:finite}
\status{checked by JS for $\Sigma = \sigma^2 I$; general $\Sigma$ not checked}
If $s(x)$ takes at most $K$ distinct values on $\mathcal{X}$, then
$R = \sqrt{2\ln(K/\epsilon)}$ guarantees
$\mathbb{P}\bigl(f(x^\star;\theta^\star) > b\bigr) \le \epsilon$.
\end{corollary}

\begin{corollary}[Linear models]\label{cor:subspace}
\status{checked by JS for $\Sigma = \sigma^2 I$; general $\Sigma$ not checked}
Suppose $s(x)$ lies in a subspace $S \subseteq \mathbb{R}^n$ of dimension $d$
for all $x \in \mathcal{X}$, and $\xi \sim N(0, I)$, that is,
$\varepsilon \sim N(0,\Sigma)$. Then $R = \sqrt{\chi^2_{d,1-\epsilon}}$
guarantees $\mathbb{P}\bigl(f(x^\star;\theta^\star) > b\bigr) \le \epsilon$,
where $\chi^2_{d,1-\epsilon}$ is the $(1-\epsilon)$-quantile of the chi-square
distribution with $d$ degrees of freedom.
\end{corollary}

\begin{corollary}[Multiple learned constraints]\label{cor:multi}
\status{checked by JS for $\Sigma = \sigma^2 I$; general $\Sigma$ not checked}
Consider $m$ learned constraints, each satisfying
Assumptions~\ref{ass:noise} and~\ref{ass:affine} with its own $s_j$,
$\Sigma_j$ and radius $R_j$. If each $R_j$ is chosen as in
Corollary~\ref{cor:finite} or~\ref{cor:subspace} at level $\epsilon/m$, then
the probability that $x^\star$ violates at least one clean-label constraint is
at most $\epsilon$. No assumption is required on the dependence of the noise
across constraints.
\end{corollary}

\begin{remark}
\status{checked by JS for $\Sigma = \sigma^2 I$}
(i) For ridge regression, $s(x) = Z(Z^\top Z + \lambda I)^{-1}\phi(x)$, so $S$
is the column space of the training matrix $Z$ and $d = \mathrm{rank}(Z)$.
(ii) A lower-bound constraint $f(x;\theta) \ge b$ is handled symmetrically,
with robust counterpart $f(x;\hat\theta) - R\|s(x)\|_\Sigma \ge b$ and the same
radii.
(iii) The guarantee concerns the propagation of label noise only; the
approximation error $h^\star - f(\cdot\,;\theta^\star)$ is not covered.
(iv) For models whose structure is selected from data (active sets, splits,
hidden layers), Assumption~\ref{ass:affine} requires the structure to be
chosen independently of $\varepsilon$, for instance by sample splitting.
\end{remark}

\paragraph{The simulated radius.}
Corollaries~\ref{cor:finite} and~\ref{cor:subspace} bound $W$ without using
$\mathcal{X}$: they protect against every direction among the vectors
$\Sigma^{1/2}s(x)$, whether or not a decision in $\mathcal{X}$ reaches it. When
the distribution of $\xi$ is known, the radius can instead be read off
simulated copies of $W$, and an order statistic of the copies is valid with no
sampling-error term. Write $W = w(\xi)$, where
\[
  w(z) = \sup_{x \in \mathcal{X},\, \|s(x)\|_\Sigma \neq 0}
    \frac{-(\Sigma^{1/2}s(x))^\top z}{\|s(x)\|_\Sigma},
  \qquad z \in \mathbb{R}^n ,
\]
assumed measurable. Under Assumption~\ref{ass:affine}, $s(\cdot)$ does not
depend on $\varepsilon$, so $w$ is a fixed function, known once the model's
structure is. Evaluating $w(z)$ is a maximization over $\mathcal{X}$, a
mixed-integer program when $\mathcal{X}$ is mixed-integer.

\begin{proposition}[Simulated radius]\label{prop:simulated}
\status{not checked: drafted 2026-11-15}
Suppose Assumptions~\ref{ass:noise} and~\ref{ass:affine} hold and the
distribution of $\xi$ is known. Draw $\xi^{(1)},\dots,\xi^{(N)}$ from it,
independently of each other and of $\varepsilon$, and let
$\overline W_1,\dots,\overline W_N$ be any numbers with
$\overline W_i \ge w(\xi^{(i)})$. For $\epsilon \in (0,1)$ and
$N \ge 1/\epsilon - 1$, let $k = \lceil (N+1)(1-\epsilon) \rceil$ and let
$\hat R$ be the $k$-th smallest of $\overline W_1,\dots,\overline W_N$. Then
any $x^\star$ satisfying \eqref{eq:rc} at $R = \hat R$, possibly chosen as a
function of the data and of the simulation, satisfies
$\mathbb{P}\bigl(f(x^\star;\theta^\star) > b\bigr) \le \epsilon$, where the
probability is over the noise and the simulation jointly.
\end{proposition}

\begin{remark}[Computing the radius]\label{rem:simulated}
\status{not checked: drafted 2026-11-15}
(i) Each $\overline W_i$ must be an upper bound on $w(\xi^{(i)})$. The value
of a feasible point of the maximization, or the largest standardized error
over a sample of decisions, is a lower bound, and an order statistic of lower
bounds can fall below the $(1-\epsilon)$-quantile of $W$, which makes the
radius too small. A solver's dual bound is an upper bound at any time limit,
so stopping a maximization early makes the radius larger, never invalid.
(ii) The proposition needs the distribution of $\xi$, not only a variance
proxy: the argument requires the simulated copies to have the law of the
true $W$. The experiments use $\xi \sim N(0, I)$. An estimated $\Sigma$ is
not covered yet. \tbd{T2}
\end{remark}

Corollary~\ref{cor:multi} holds with unequal levels $\epsilon_j$ summing to
$\epsilon$, and with any radius valid at level $\epsilon_j$ for constraint $j$
alone, including the simulated radius of
Proposition~\ref{prop:simulated}: its proof uses only
$\mathbb{P}(W_j > R_j) \le \epsilon_j$. \status{not checked: drafted
2026-11-15}
) in Section 3. Uses the
% assumption, definition, theorem, corollary and remark environments and the
% \status macro defined there. Citation keys are those of paper_a.bib.
%
% Proofs other than Theorem 1's are in paper_a_proofs.tex (the appendix).
% Results held until probe P1 decides the chemotherapy design (joint radius,
% certified level, simulated level) are in paper_a_guarantee_held.tex, which
% is not \input.
%
% \status tags. JS checked by hand the version committed in 40acea5: Theorem 1,
% Corollaries 1-3 and the closing remark, for Sigma = sigma^2 I. The general
% Sigma, the remark on the choice of norm and everything from "The simulated
% radius" on came later and are not checked.

\begin{outline}
Tags: JS checked by hand the $\Sigma = \sigma^2 I$ version of
Theorem~\ref{thm:selection}, Corollaries~\ref{cor:finite}--\ref{cor:multi}
and the closing remark (commit \texttt{40acea5}). The general $\Sigma$, the
remark on the choice of norm and everything after the closing remark were
added later and are not checked.
\end{outline}

\paragraph{Setting.}
We observe training data $(x_i, y_i)_{i=1}^n$ with $y = y^\star + \varepsilon$,
where $y^\star_i = h^\star(x_i)$ are the noise-free labels. The learned
constraint is $f(x;\theta(y)) \le b$ for $x \in \mathcal{X}$, where $\theta(y)$
denotes the model fitted on labels $y$.

\begin{assumption}[Label noise]\label{ass:noise}
\status{checked by JS for $\Sigma = \sigma^2 I$; general $\Sigma$ not checked}
The noise is $\varepsilon = \Sigma^{1/2}\xi$, where $\Sigma \succeq 0$ is known
and $\xi_1,\dots,\xi_n$ are independent, zero-mean and sub-Gaussian with
variance proxy $1$. We write $\|v\|_\Sigma = \|\Sigma^{1/2}v\|_2$.
\end{assumption}

With $\Sigma = \sigma^2 I$ the noise terms are independent with variance proxy
$\sigma^2$ and $\|v\|_\Sigma = \sigma\|v\|_2$. A diagonal $\Sigma$ allows
unequal noise levels, and a full one correlated noise.

\begin{assumption}[Affine predictor]\label{ass:affine}
For all $x \in \mathcal{X}$, $f(x;\theta(y)) = s(x)^\top y + c(x)$, where
$s(x) \in \mathbb{R}^n$ and $c(x)$ do not depend on $\varepsilon$.
\end{assumption}

\begin{definition}[Clean-label model]\label{def:clean}
Let $\hat\theta = \theta(y)$ and $\theta^\star = \theta(y^\star)$. The
clean-label model is $f(\cdot\,;\theta^\star)$.
\end{definition}

Since $s(x)$ and $c(x)$ do not depend on $\varepsilon$,
Assumption~\ref{ass:affine} applies to $y^\star = y - \varepsilon$ as well as to
$y$, and subtracting the two gives, for all $x \in \mathcal{X}$,
\begin{equation}\label{eq:identity}
  f(x;\theta^\star) = f(x;\hat\theta) - s(x)^\top \varepsilon .
\end{equation}

For a radius $R > 0$ fixed independently of $\varepsilon$, the robust
counterpart of the learned constraint is
\begin{equation}\label{eq:rc}
  f(x;\hat\theta) + R\,\|s(x)\|_\Sigma \le b .
\end{equation}
The radius $R$ is in units of the noise's scale: with $\Sigma = \sigma^2 I$ the
margin is $R\sigma\|s(x)\|_2$.

\begin{theorem}\label{thm:selection}
\status{checked by JS for $\Sigma = \sigma^2 I$; general $\Sigma$ not checked}
Suppose Assumptions~\ref{ass:noise} and~\ref{ass:affine} hold, and let $x^\star$
be any solution satisfying \eqref{eq:rc}, possibly chosen as a function of the
data. Then
\[
  \mathbb{P}\bigl(f(x^\star;\theta^\star) > b\bigr) \le \mathbb{P}(W > R),
  \qquad
  W = \sup_{x \in \mathcal{X},\, \|s(x)\|_\Sigma \neq 0}
      \frac{-s(x)^\top \varepsilon}{\|s(x)\|_\Sigma} .
\]
\end{theorem}

\begin{proof}
Suppose $f(x^\star;\theta^\star) > b$. By \eqref{eq:identity} and
\eqref{eq:rc},
\[
  -s(x^\star)^\top \varepsilon
  = f(x^\star;\theta^\star) - f(x^\star;\hat\theta)
  > b - f(x^\star;\hat\theta)
  \ge R\,\|s(x^\star)\|_\Sigma .
\]
The left side equals $-(\Sigma^{1/2}s(x^\star))^\top\xi$ and is nonzero, so
$\|s(x^\star)\|_\Sigma \neq 0$ and $W > R$.
\end{proof}

\begin{remark}[Choice of norm]\label{rem:norm}
\status{not checked: added after JS's check}
The proof uses only the inequality
$-s(x^\star)^\top\varepsilon > R\,\|s(x^\star)\|_\Sigma$, so
Theorem~\ref{thm:selection} holds verbatim for any norm $N$ on $\mathbb{R}^n$
in place of $\|\cdot\|_\Sigma$, both in \eqref{eq:rc} and in $W$. The norm
enters only through the distribution of $W$. Under $\|\cdot\|_\Sigma$ the
standardized error at $x$ is
$-s(x)^\top\varepsilon/\|s(x)\|_\Sigma = -u(x)^\top\xi$ with $u(x)$ a unit
vector, so it has the same tail at every $x$; Corollaries~\ref{cor:finite}
and~\ref{cor:subspace} use exactly this. Under another norm the standardized
error at $x$ is that one multiplied by
$\rho(x) = \|s(x)\|_\Sigma / N(s(x))$. Bounding the corresponding $W_N$
through $W$, as $W_N \le \bar\rho\max(W,0)$ with
$\bar\rho = \sup_{x\in\mathcal{X}} \rho(x)$, turns a radius $R$ from these
corollaries into $\bar\rho R$, and the margin
$\bar\rho R\,N(s(x)) \ge R\,\|s(x)\|_\Sigma$ is at least as large at every
$x$. For example, noise bounded as $|\varepsilon_i| \le c$ gives the box
counterpart $f(x;\hat\theta) + c\,\|s(x)\|_1 \le b$, which holds with
probability one; but $\|s(x)\|_1$ can be up to $\sqrt{n}$ times
$\|s(x)\|_2$.
\end{remark}

\begin{corollary}[Tree ensembles]\label{cor:finite}
\status{checked by JS for $\Sigma = \sigma^2 I$; general $\Sigma$ not checked}
If $s(x)$ takes at most $K$ distinct values on $\mathcal{X}$, then
$R = \sqrt{2\ln(K/\epsilon)}$ guarantees
$\mathbb{P}\bigl(f(x^\star;\theta^\star) > b\bigr) \le \epsilon$.
\end{corollary}

\begin{corollary}[Linear models]\label{cor:subspace}
\status{checked by JS for $\Sigma = \sigma^2 I$; general $\Sigma$ not checked}
Suppose $s(x)$ lies in a subspace $S \subseteq \mathbb{R}^n$ of dimension $d$
for all $x \in \mathcal{X}$, and $\xi \sim N(0, I)$, that is,
$\varepsilon \sim N(0,\Sigma)$. Then $R = \sqrt{\chi^2_{d,1-\epsilon}}$
guarantees $\mathbb{P}\bigl(f(x^\star;\theta^\star) > b\bigr) \le \epsilon$,
where $\chi^2_{d,1-\epsilon}$ is the $(1-\epsilon)$-quantile of the chi-square
distribution with $d$ degrees of freedom.
\end{corollary}

\begin{corollary}[Multiple learned constraints]\label{cor:multi}
\status{checked by JS for $\Sigma = \sigma^2 I$; general $\Sigma$ not checked}
Consider $m$ learned constraints, each satisfying
Assumptions~\ref{ass:noise} and~\ref{ass:affine} with its own $s_j$,
$\Sigma_j$ and radius $R_j$. If each $R_j$ is chosen as in
Corollary~\ref{cor:finite} or~\ref{cor:subspace} at level $\epsilon/m$, then
the probability that $x^\star$ violates at least one clean-label constraint is
at most $\epsilon$. No assumption is required on the dependence of the noise
across constraints.
\end{corollary}

\begin{remark}
\status{checked by JS for $\Sigma = \sigma^2 I$}
(i) For ridge regression, $s(x) = Z(Z^\top Z + \lambda I)^{-1}\phi(x)$, so $S$
is the column space of the training matrix $Z$ and $d = \mathrm{rank}(Z)$.
(ii) A lower-bound constraint $f(x;\theta) \ge b$ is handled symmetrically,
with robust counterpart $f(x;\hat\theta) - R\|s(x)\|_\Sigma \ge b$ and the same
radii.
(iii) The guarantee concerns the propagation of label noise only; the
approximation error $h^\star - f(\cdot\,;\theta^\star)$ is not covered.
(iv) For models whose structure is selected from data (active sets, splits,
hidden layers), Assumption~\ref{ass:affine} requires the structure to be
chosen independently of $\varepsilon$, for instance by sample splitting.
\end{remark}

\paragraph{The simulated radius.}
Corollaries~\ref{cor:finite} and~\ref{cor:subspace} bound $W$ without using
$\mathcal{X}$: they protect against every direction among the vectors
$\Sigma^{1/2}s(x)$, whether or not a decision in $\mathcal{X}$ reaches it. When
the distribution of $\xi$ is known, the radius can instead be read off
simulated copies of $W$, and an order statistic of the copies is valid with no
sampling-error term. Write $W = w(\xi)$, where
\[
  w(z) = \sup_{x \in \mathcal{X},\, \|s(x)\|_\Sigma \neq 0}
    \frac{-(\Sigma^{1/2}s(x))^\top z}{\|s(x)\|_\Sigma},
  \qquad z \in \mathbb{R}^n ,
\]
assumed measurable. Under Assumption~\ref{ass:affine}, $s(\cdot)$ does not
depend on $\varepsilon$, so $w$ is a fixed function, known once the model's
structure is. Evaluating $w(z)$ is a maximization over $\mathcal{X}$, a
mixed-integer program when $\mathcal{X}$ is mixed-integer.

\begin{proposition}[Simulated radius]\label{prop:simulated}
\status{not checked: drafted 2026-11-15}
Suppose Assumptions~\ref{ass:noise} and~\ref{ass:affine} hold and the
distribution of $\xi$ is known. Draw $\xi^{(1)},\dots,\xi^{(N)}$ from it,
independently of each other and of $\varepsilon$, and let
$\overline W_1,\dots,\overline W_N$ be any numbers with
$\overline W_i \ge w(\xi^{(i)})$. For $\epsilon \in (0,1)$ and
$N \ge 1/\epsilon - 1$, let $k = \lceil (N+1)(1-\epsilon) \rceil$ and let
$\hat R$ be the $k$-th smallest of $\overline W_1,\dots,\overline W_N$. Then
any $x^\star$ satisfying \eqref{eq:rc} at $R = \hat R$, possibly chosen as a
function of the data and of the simulation, satisfies
$\mathbb{P}\bigl(f(x^\star;\theta^\star) > b\bigr) \le \epsilon$, where the
probability is over the noise and the simulation jointly.
\end{proposition}

\begin{remark}[Computing the radius]\label{rem:simulated}
\status{not checked: drafted 2026-11-15}
(i) Each $\overline W_i$ must be an upper bound on $w(\xi^{(i)})$. The value
of a feasible point of the maximization, or the largest standardized error
over a sample of decisions, is a lower bound, and an order statistic of lower
bounds can fall below the $(1-\epsilon)$-quantile of $W$, which makes the
radius too small. A solver's dual bound is an upper bound at any time limit,
so stopping a maximization early makes the radius larger, never invalid.
(ii) The proposition needs the distribution of $\xi$, not only a variance
proxy: the argument requires the simulated copies to have the law of the
true $W$. The experiments use $\xi \sim N(0, I)$. An estimated $\Sigma$ is
not covered yet. \tbd{T2}
\end{remark}

Corollary~\ref{cor:multi} holds with unequal levels $\epsilon_j$ summing to
$\epsilon$, and with any radius valid at level $\epsilon_j$ for constraint $j$
alone, including the simulated radius of
Proposition~\ref{prop:simulated}: its proof uses only
$\mathbb{P}(W_j > R_j) \le \epsilon_j$. \status{not checked: drafted
2026-11-15}
